
\bigskip

%Les parties A et B peuvent être traitées de façon indépendante.
%
%\medskip
%
%Au cours de la fabrication d'une paire de lunettes, la paire de verres doit subir deux traitements notés T1 et T2.
%
%\bigskip

\textbf{Partie A}

\medskip

%On prélève au hasard une paire de verres dans la production.
%
%On désigne par $A$ l'évènement : \og la paire de verres présente un défaut pour le traitement T1 \fg.
%
%On désigne par $B$ l'évènement : \og la paire de verres présente un défaut pour le traitement T2 \fg.
%
%On note respectivement $\overline{A}$ et $\overline{B}$ les évènements contraires de $A$ et $B$.
%
%Une étude a montré que :
%
%\begin{itemize}
%\item la probabilité qu'une paire de verres présente un défaut pour le traitement T1 notée $P(A)$ est égale à 0,1.
%\item la probabilité qu'une paire de verres présente un défaut pour le traitement T2 notée $P(B)$ est égale à 0,2.
%\item la probabilité qu'une paire de verres ne présente aucun des deux défauts est 0,75.
%\end{itemize}

\medskip

\begin{enumerate}
\item Recopier et compléter le tableau suivant avec les probabilités correspondantes.

\begin{center}
\begin{tabularx}{0.45\linewidth}{|*{4}{>{\centering \arraybackslash}X|}}\hline
				&$A$ 	&$\overline{A}$	&Total\\ \hline
\black $B$				&\blue 0,05		&\blue 0,15				&\black 0,2\\ \hline
\black $\overline{B}$	&\blue 0,05		&\black 0,75				&\blue 0,8\\ \hline
\black Total			&\black 0,1		&\blue 0,9				&\black 1\\ \hline
\end{tabularx}
\end{center}

\item
	\begin{enumerate}
		\item %Déterminer, en justifiant la réponse, la probabilité qu'une paire de verres, prélevée au hasard dans la production, présente un défaut pour au moins un des deux traitements T1 ou T2.
La probabilité qu'une paire de verres, prélevée au hasard dans la production, présente un défaut pour au moins un des deux traitements T1 ou T2 est l'évènement contraire de l'évènement \og une paire de verres ne présente aucun des deux défauts \fg, donc sa probabilité est égale $1 - 0,75 = 0,25$. Ou encore $P(A \cup B) = 1 - \left(\overline{A} \cap \overline{B} \right)$.
		\item %Donner la probabilité qu'une paire de verres, prélevée au hasard dans la production, présente deux défauts, un pour chaque traitement T1 et T2.
		
$\bullet~~$On a $P\left(\overline{A}\right) = 1 - P(A) = 1 - 0,1 = 0,9$ ;

$\bullet~~$On a $P\left(\overline{B}\right) = 1 - P(B) = 1 - 0,2 = 0,8$.

Donc $P\left(A \cap \overline{B}\right) = 0,8 - 0,75 = 0,05$ et 

$P\left(B \cap \overline{A}\right) = 0,9 - 0,75 = 0,15$.

Donc la probabilité qu'une paire de verres, prélevée au hasard dans la production, présente deux  défauts  est égale à :

$P(A \cap B) = 1 - P\left(A \cap \overline{B}\right)  -  P\left(B \cap \overline{A}\right) - P\left(\overline{A} \cap \overline{B} \right) = 1 - 0,05 - 0,15  - 0,75 = 0,05$.
		\item %Les évènements $A$ et $B$ sont-ils indépendants? Justifier la réponse.
$P(A) \times P(B) = 0,1 \times 0,2 = 0,02$ et $P(A \cap B) = 0,05$.

On a donc $P(A \cap B) \ne P(A) \times P(B)$ : les évènements $A$ et $B$ ne sont pas indépendants.
	\end{enumerate}	
\item %Calculer la probabilité qu'une paire de verres, prélevée au hasard dans la production, présente un défaut pour un seul des deux traitements.
On a $P\left(A \cap \overline{B}\right) + P\left(B \cap \overline{A}\right)= 0,05 + 0,15 = 0,2$
\item %Calculer la probabilité qu'une paire de verres, prélevée au hasard dans la production, présente un défaut pour le traitement T2, sachant que cette paire de verres présente un défaut pour le traitement T1.
On a $P_{A}\left(B \right) = \dfrac{P(A \cap B)}{P(A)} = \dfrac{0,05}{0,1} = \dfrac12$.
\end{enumerate}

\bigskip

\textbf{Partie B}

\medskip

%On prélève, au hasard, un échantillon de $50$ paires de verres dans la production. On suppose que la production est suffisamment importante pour assimiler ce prélèvement à un tirage avec remise. 
%
%On note $X$ la variable aléatoire qui, à chaque échantillon de ce type, associe le nombre de paires de verres qui présentent le défaut pour le traitement T1.
%
%\medskip

\begin{enumerate}
\item %Justifier que la variable aléatoire $X$ suit une loi binomiale et préciser les paramètres de cette loi.
La production est suffisamment importante pour que la probabilité d'avoir le déf	ut $T1$ est égale à 0,1. Comme il y a 50 tirages indépendants, la variable aléatoire $X$ suit une loi binomiale de paramètres $n = 50$ et $p = 0,1$.
\item %Donner l'expression permettant de calculer la probabilité d'avoir, dans un tel échantillon, exactement $10$ paires de verres qui présentent ce défaut. 
On sait que $P(X = 10) = \binom{50}{10}\times 0,1^{10} \times (1 - 0,1)^{50 - 10}$.

La calculatrice donne $P(X = 10)  \approx 0,015$ au millième près.

%Effectuer ce calcul et arrondir le résultat à $10^{-3}$.
\item %En moyenne, combien de paires de verres ayant ce défaut peut-on trouver dans un échantillon de $50$ paires ?
La moyenne est l'espérance de la variable $X$ et on sait que $E(X) = n \times p = 50 \times 0,1 = 5$.
\end{enumerate}
